\needspace{0.25\textheight}
\section{Уравнения Эйлера--Пуанкаре с диссипацией}\label{sec:ep}

\subsection{Вариационный вывод}

Зафиксируем гладкую кривую $g(t) \in \mathrm{SO}(3)$, $t \in [0,T]$,
и следом за \cite{arnold,hmr,marsdenratiu} положим
\[
  \hat{\om}(t) := g(t)^{-1}\dot g(t) \in \so,
  \qquad
  \Om(t) := \bbI\,\om(t) \in \so^{*} \cong \R^3 .
\]
Редуцированный лагранжиан~--- кинетическая энергия вихревой трубки:
$\ell(\om) = \tfrac12\langle \om, \bbI\om\rangle$; тогда
$\delta\ell/\delta\om = \Om$ (отождествление $\so^{*} \cong \so$~---
метрикой Киллинга, лемма~\ref{lem:killing}).

\begin{lemma}[об усечённых вариациях]\label{lem:variation}
Пусть $g_\varepsilon$~--- гладкая вариация кривой $g$ с закреплёнными
концами: $g_\varepsilon(0) \equiv g(0)$, $g_\varepsilon(T) \equiv g(T)$.
Положим $\hat\eta := g^{-1}\delta g$, где
$\delta g = \left.\partial_\varepsilon g_\varepsilon\right|_{\varepsilon=0}$.
Тогда $\eta(0) = \eta(T) = 0$ и
\begin{equation}\label{eq:varomega}
  \delta\hat{\om} = \dot{\hat\eta} + [\hat\om, \hat\eta] .
\end{equation}
\end{lemma}

\begin{proof}
Тождество $g_\varepsilon g_\varepsilon^{-1} \equiv \mathrm{Id}$ даёт
при дифференцировании по $\varepsilon$:
$\delta(g^{-1}) = -g^{-1}\,\delta g\,g^{-1}$. Далее, из
$\hat\om = g^{-1}\dot g$:
\[
  \delta\hat\om
  = \delta(g^{-1})\,\dot g + g^{-1}\,\delta(\dot g)
  = -g^{-1}\delta g\,g^{-1}\dot g + g^{-1}\partial_t(\delta g).
\]
Так как $\delta g = g\hat\eta$, то
$\partial_t(\delta g) = \dot g\,\hat\eta + g\,\dot{\hat\eta}$, и потому
\[
  \delta\hat\om = -\hat\eta\hat\om + \hat\om\hat\eta + \dot{\hat\eta}
  = \dot{\hat\eta} + [\hat\om,\hat\eta].
  \qedhere
\]
\end{proof}

\begin{theorem}[Эйлер--Пуанкаре]\label{th:ep}
Кривая $g(\cdot)$~--- критическая точка функционала
$\int_0^T \ell(\om)\,dt$ относительно вариаций с закреплёнными концами
тогда и только тогда, когда
\begin{equation}\label{eq:ep}
  \dot\Om = \adst_{\om}\,\Om = \Om \times \om = \Om \times \bbI^{-1}\Om .
\end{equation}
\end{theorem}

\begin{proof}
По лемме~\ref{lem:variation} и правилу дифференцирования композиции,
\[
  \delta \int_0^T \ell(\om)\,dt
  = \int_0^T \langle \Om,\, \dot\eta + \om\times\eta\rangle\,dt
  = -\int_0^T \langle \dot\Om,\,\eta\rangle\,dt
    + \int_0^T \langle \Om,\,\om\times\eta\rangle\,dt ,
\]
где граничный член $[\langle\Om,\eta\rangle]_0^T$ обнуляется, ибо
$\eta(0) = \eta(T) = 0$. Второй интеграл перепишем через коалгебраическое
действие: по определению
$\langle \adst_{\om}\Om, \eta\rangle := \langle \Om, [\om,\eta]\rangle
= \Om\cdot(\om\times\eta)$; цикличность смешанного произведения даёт
\begin{equation}\label{eq:adstar}
  \Om\cdot(\om\times\eta) = \eta\cdot(\Om\times\om)
  \qquad\Longrightarrow\qquad
  \adst_{\om}\Om = \Om\times\om .
\end{equation}
Итого $\delta S = \int_0^T \langle -\dot\Om + \Om\times\om,\;\eta\rangle dt$,
и обращение в нуль для всех $\eta$ равносильно \eqref{eq:ep}.
\end{proof}

\subsection{Диссипативное расширение и постановка задачи}

Вязкое слагаемое уравнения \eqref{eq:ns-vort} в редуцированной модели
представляется линейным оператором затухания. Далее всюду
\begin{equation}\label{eq:epd}
  \boxed{\;
  \dot\Om = \underbrace{\Om\times(\bbI^{-1}\Om)}_{\text{поворот фазы}}
  \;-\; \nu\,\bbL\,\Om,\qquad
  \Om(0) = \Om_0,}
\end{equation}
при гипотезах
\begin{itemize}[leftmargin=2.2em, itemsep=1pt]
  \item[\textup{(H1)}] $\bbL = \bbL^{\mathsf T} \succeq 0$~--- всюду;
  \item[\textup{(H2)}] $[\bbL, \bbI] = 0$ (эквивалентно: $\bbL$ диагонален
        в главном базисе $\bbI$)~--- там, где указано явно.
\end{itemize}
Слагаемое $\nu\bbL\Om$~--- структурный аналог $\nu\Delta\bwn$: оператор
$\bbL$ играет роль положительного «символа» $-\Delta$, что уточняется
балансом энстрофии в \S\ref{sec:cascade}. В покомпонентной записи
\eqref{eq:epd} принимает вид
\begin{equation}\label{eq:components}
  \begin{aligned}
    \dot\Omega_1 &= \Bigl(\tfrac1{I_3}-\tfrac1{I_2}\Bigr)\Omega_2\Omega_3 - \nu L_{11}\Omega_1,\\
    \dot\Omega_2 &= \Bigl(\tfrac1{I_1}-\tfrac1{I_3}\Bigr)\Omega_3\Omega_1 - \nu L_{22}\Omega_2,\\
    \dot\Omega_3 &= \Bigl(\tfrac1{I_2}-\tfrac1{I_1}\Bigr)\Omega_1\Omega_2 - \nu L_{33}\Omega_3,
  \end{aligned}
\end{equation}
где в последней строке использована (H2); при $\nu = 0$ это классические
уравнения Эйлера для твёрдого тела \cite{euler,arnold}.

\subsection{Гамильтонова структура: скобка Ли--Пуансона}

\begin{proposition}\label{prop:poisson}
На $C^\infty(\R^3)$ скобка
\begin{equation}\label{eq:lp}
  \{F,G\}(\Om) := -\bigl\langle \Om,\; \nabla F(\Om)\times\nabla G(\Om)\bigr\rangle
\end{equation}
корректно определена, билинейна, кососимметрична и удовлетворяет
тождеству Якоби; функция
\[
  C(\Om) = \tfrac12\|\Om\|^2
\]
является Казимиром скобки: $\{C,F\} \equiv 0$ для всех $F$. Система
\eqref{eq:ep} есть гамильтонова система $\dot f = \{f,\mathcal H\}$
с гамильтонианом $\mathcal H(\Om) = \tfrac12\langle\Om,\bbI^{-1}\Om\rangle$.
\end{proposition}

\begin{proof}
Билинейность и кососимметрия очевидны. \emph{Тождество Якоби.}
Обозначим $x = \nabla X(\Om)$, $y = \nabla Y(\Om)$, $z = \nabla Z(\Om)$,
$X'' = \nabla^2X(\Om)$ и т.д. Дифференцирование произведения в
\eqref{eq:lp} даёт (используется симметрия гессианов):
\begin{equation}\label{eq:gradbracket}
  \nabla\{X,Y\}(\Om) = -(x\times y) - X''\,(y\times\Om) - Y''\,(\Om\times x).
\end{equation}
Действительно, $\partial_i\bigl[\Om_\ell\varepsilon_{\ell mn}\partial_mX\,\partial_nY\bigr]
= \varepsilon_{imn}x_my_n + \Om_\ell\varepsilon_{\ell mn}
\bigl(X''_{im}y_n + x_mY''_{in}\bigr)$, и вторая группа сворачивается как
$\bigl[X''(y\times\Om) + Y''(\Om\times x)\bigr]_i$
(перестановка циклов смешанного произведения; симметрия $X''$, $Y''$
существенна). Подставляя \eqref{eq:gradbracket} в \eqref{eq:lp} с
внешней парой $(\,\{X,Y\}, Z\,)$ и применяя формулу Лагранжа:
\begin{equation}\label{eq:closedform}
  \{\{X,Y\},Z\}(\Om)
  = (\Om\cdot y)(x\cdot z) - (\Om\cdot x)(y\cdot z)
  + (y\times\Om)^{\mathsf T} X''\,(z\times\Om)
  + (\Om\times x)^{\mathsf T} Y''\,(z\times\Om).
\end{equation}
Просуммируем \eqref{eq:closedform} и её циклические перестановки
$(X,Y,Z) \to (Y,Z,X) \to (Z,X,Y)$. Бесгессиановая часть обнуляется:
слагаемое $(\Om\cdot v)(w\cdot u)$ входит в сумму дважды~---
с плюсом и с минусом~--- по каждому из трёх векторов. Гессиановые
члены сокращаются парами: $X''$-члены входят в
$\{\{X,Y\},Z\}$ и $\{\{Z,X\},Y\}$ как
\[
  (y\times\Om)^{\mathsf T}X''(z\times\Om)
  \;+\;(\Om\times z)^{\mathsf T}X''(y\times\Om)
  = p - p^{\mathsf T} = 0,
\]
где $p := (y\times\Om)^{\mathsf T}X''(z\times\Om)$~--- скаляр, равный
своему транспонированному виду
$p^{\mathsf T} = (z\times\Om)^{\mathsf T}X{}''^{\mathsf T}(y\times\Om)
= (z\times\Om)^{\mathsf T}X''(y\times\Om) = -(\Om\times z)^{\mathsf T}X''(y\times\Om)$,
и снова использована симметрия $X''$. Аналогично для $Y''$ и $Z''$.
Итого циклическая сумма равна нулю~--- тождество Якоби.

\emph{Казимир.} $\{C,F\} = -\langle \Om, \Om\times\nabla F\rangle
= -\det[\Om,\Om,\nabla F] = 0$.

\emph{Гамильтоновость.} Для $f \in C^\infty$ вдоль \eqref{eq:ep}:
$\dot f = \langle\nabla f, \Om\times\bbI^{-1}\Om\rangle$; с другой
стороны, $\{f,\mathcal H\} = -\langle\Om, \nabla f\times\bbI^{-1}\Om\rangle
= \langle \nabla f, \Om\times\bbI^{-1}\Om\rangle$ (цикличность).
Совпадение даёт $\dot f = \{f,\mathcal H\}$; при $\nabla f = e_k$
получаем покомпонентно \eqref{eq:ep}.
\end{proof}

\begin{remark}
Формула \eqref{eq:closedform}~--- замкнутое выражение двойной скобки
Ли--Пуансона на $\so^{*}$; стандартные доказательства (см.
\cite{marsdenratiu,arnoldkhesin}) сводят тождество Якоби к
структурным константам алгебры. Приведённое выше прямое доказательство
дополнительно даёт явный вид сокращения гессиановых членов
($p - p^{\mathsf T} = 0$).
\end{remark}

\subsection{Корректность задачи Коши}

\begin{theorem}[глобальная корректность]\label{th:wp}
При любых $\Om_0 \in \R^3$, $\nu \ge 0$ и $\bbL$, удовлетворяющей
\textup{(H1)}, задача Коши \eqref{eq:epd} имеет единственное
$C^\infty$-решение на $[0,\infty)$.
\end{theorem}

\begin{proof}
Правая часть $f(\Om) = \Om\times\bbI^{-1}\Om - \nu\bbL\Om$ гладка.
Оценка на шаре $B(0,R)$:
\[
  \|f(\Om) - f(\tilde\Om)\|
  \le \|\bbI^{-1}\|\,\bigl(\|\Om\| + \|\tilde\Om\|\bigr)\,\|\Om-\tilde\Om\|
     + \nu\|\bbL\|\,\|\Om-\tilde\Om\|
  \le \bigl(2\|\bbI^{-1}\|R + \nu\|\bbL\|\bigr)\|\Om-\tilde\Om\|,
\]
где использовано $\|a\times p - b\times q\|
\le \|(a-b)\times p\| + \|b\times(p-q)\| \le \|a-b\|\|p\| + \|b\|\|p-q\|$
и $\|\bbI^{-1}v\| \le \|\bbI^{-1}\|\,\|v\|$. По теореме
Пикара--Линделёфа существует единственное локальное решение.
Априорная оценка: вдоль решения
\[
  \frac{d}{dt}\|\Om\|^2 = 2\bigl\langle \Om,\Om\times\bbI^{-1}\Om\bigr\rangle
  - 2\nu\,\langle\Om,\bbL\Om\rangle = -2\nu\,\langle\Om,\bbL\Om\rangle \le 0
\]
(тройное произведение с двумя равными сомножителями равно нулю;
$\bbL \succeq 0$ по \textup{(H1)}). Значит $\|\Om(t)\| \le \|\Om_0\|$,
решение остаётся в компактном шаре, и критерий продолжения
\cite[гл.~II]{hartman} даёт продолжение на всё $[0,\infty)$.
\end{proof}
